\clearpage\thispagestyle{empty}
\pdfbookmark[0]{Quellen zu Pallas (14. Januar 1811)}{pallas1811jan-sources}
\begin{center}\textbf{Quellen zu Pallas (14. Januar 1811)}\end{center}
\noindent Diese KI-Übertragung ins heutige Deutsch folgt dem Bericht in \href{https://gdz.sub.uni-goettingen.de/id/PPN236007467}{Gauß’ \textit{Werken}, Bd.~VI} (Göttingen, 1874), S.~322--324. Der Bericht erschien unter der Überschrift Göttingen in der Zeitschrift \textit{Göttingische gelehrte Anzeigen}, Stück 8 und 9 vom 14. Januar 1811, S.~\href{https://api.digitale-sammlungen.de/iiif/image/v2/bsb10538679_00109/full/1600,/0/default.jpg}{73}, \href{https://api.digitale-sammlungen.de/iiif/image/v2/bsb10538679_00110/full/1600,/0/default.jpg}{74}, \href{https://api.digitale-sammlungen.de/iiif/image/v2/bsb10538679_00111/full/1600,/0/default.jpg}{75}, \href{https://api.digitale-sammlungen.de/iiif/image/v2/bsb10538679_00112/full/1600,/0/default.jpg}{76} und \href{https://api.digitale-sammlungen.de/iiif/image/v2/bsb10538679_00113/full/1600,/0/default.jpg}{77}. Das digitalisierte Exemplar stammt aus der Bayerischen Staatsbibliothek.
\par\smallskip\noindent In der Korrekturtabelle der Werkausgabe steht vor dem Deklinationswert 13 für den 27. Dezember ein Zeichen mit einem senkrechten Strich und einem rechten Arm. Im Erstdruck steht ein vollständiges Pluszeichen. Das Zeichen der Werkausgabe ist als bearbeitbare Kontur wiedergegeben. Die geschätzte südliche Deklination bei der Seeberger Beobachtung lautet in der Werkausgabe $20^\circ\;58'$, auf S.~76 der Zeitschrift $21^\circ\;58'$. Die Ephemeride vom Oktober steht in der Zeitschrift \href{https://archive.org/download/bub_gb_YCI4AAAAMAAJ/page/n403.jpg}{\textit{Monatliche Correspondenz} vom Oktober 1810, S.~400--403}. KI-Übertragung und KI-Satz.
\clearpage
